\documentclass[11pt]{article}
\usepackage[T1]{fontenc}\usepackage[utf8]{inputenc}
\usepackage{lmodern,amsmath,amssymb,amsthm,mathtools,microtype,xcolor}
\usepackage[margin=1in]{geometry}
\usepackage[colorlinks=true,linkcolor=blue!50!black,urlcolor=blue!50!black,citecolor=blue!50!black]{hyperref}
\newtheorem{theorem}{Theorem}[section]\newtheorem{proposition}[theorem]{Proposition}
\newtheorem{lemma}[theorem]{Lemma}\newtheorem{corollary}[theorem]{Corollary}
\theoremstyle{remark}\newtheorem{remark}[theorem]{Remark}
\newcommand{\E}{\mathbb E}\newcommand{\Prb}{\mathbb P}\newcommand{\R}{\mathbb R}
\newcommand{\norm}[1]{\left\Vert#1\right\Vert}
\newcommand{\Lip}{\mathrm{Lip}}\newcommand{\KL}{\mathrm{KL}}
\setlength{\emergencystretch}{3em}
\hypersetup{pdftitle={Width-independent Bayesian neural-network learning on continuous input domains},pdfauthor={Kaarel Hanni; Dmitry Vaintrob; OpenAI Codex}}
\title{Width-independent Bayesian neural-network learning on continuous input domains}
\author{Kaarel H\"anni \and Dmitry Vaintrob \and OpenAI Codex\thanks{AI assistant; contributed to mathematical analysis and drafting. This research draft has not been formalized in Lean.}}
\date{October 18, 2026}
\begin{document}\maketitle
\begin{abstract}
We prove sample-complexity bounds for Bayesian tanh networks with inputs in
any compact subset of the Euclidean unit ball, including an entire sphere.
A real-valued target is implemented by a teacher of width $C$, depth $L$,
and weight bound $B$. For every learner width $N\geq C$, the sample bound
is polynomial in $d,C,B$ at fixed depth, independently of $N$.
The result holds in the strong-prior limit and at every finite-scale
exponent $\gamma\geq1$, including mean-field scaling. Gaussian weights
and a family of smooth, strongly log-concave non-Gaussian weights are covered.
Training observes clean real-valued labels and conditions on fitting them
within a positive tolerance $\rho$; testing uses a larger tolerance
$\tau>\rho$. The full posterior assigns probability greater than $2/3$
to a prediction within $\tau$ on all test inputs outside a set of
probability at most $\varepsilon$, with high probability over the sample. No restriction
on the input distribution beyond its support is needed, and no global
output-margin condition is imposed on the prior.
The continuous-domain argument uses a finite input net only inside the
proof. Bounded-energy compression handles the strong-prior and
$\gamma>1$ regimes. A bounded dual representation of an arbitrary readout
extends the mean-field covariance argument from signs to real outputs.
We also obtain exact interpolation in an iterated Gaussian strong-prior
limit, and binary classification on supports where a teacher has a
uniform positive margin. Exact equality conditioning at fixed prior scale,
unbounded teacher weights, and margin-free binary-label classification
are not covered.
\end{abstract}

\section{Introduction}
A small neural-network implementation can be a useful source of statistical
complexity even when the Bayesian learner has many more parameters. This
paper establishes such a guarantee on continuous input domains. The sample
requirement depends on the width and weight magnitude of a teacher, rather
than on the width of the learner. The dependence on depth is allowed to be
large: throughout the paper, depth is fixed.

There is a basic obstacle to directly importing a Boolean-domain model.
A continuous real-valued network whose output is bounded away from zero
everywhere on a connected domain has a constant sign. Thus conditioning a
prior to have an absolute margin at every point of a connected sphere
would exclude every nonconstant binary classifier. We instead work with
the raw independent-weight prior and condition only on the observed
training constraints. Our principal result concerns real-valued regression,
which permits nonconstant targets on the whole sphere without a margin
assumption. For example, $f(x)=\tanh(x_1)$ on $S^{d-1}$ is covered by a
width-one, depth-two teacher.

The labels in this theorem are clean, exact values of the target, but the
likelihood accepts a fit within a fixed positive tolerance. Test error
means exceeding a strictly larger tolerance. This difference between
training and test thresholds provides room for finite approximation.
Conditioning directly on exact equality of continuously distributed real
outputs at fixed prior scale is a different problem. We nevertheless obtain
an exact-interpolation corollary by taking the strong-prior limit first and
then sending the training tolerance to zero, in that specified order.

The proofs have three parts. First, a neighborhood of a cloned teacher has
enough prior mass, with an exponential rate independent of the learner
width (the mass itself may decrease with width). Second, a finite family
approximates all but a sufficiently small amount of prior mass. Third, a
union bound shows that an inaccurate member of this family is unlikely to
fit the data. For continuous domains, Lipschitz control lets us use a finite
input net whose logarithmic cardinality is polynomial in the relevant
parameters. The net is a proof device; the data distribution remains
arbitrary and continuous inputs are not discretized by the learning rule.

Several probabilistic ingredients extend the finite-input analysis
in~\cite{nongaus}. We give their statements and the necessary derivations
here, and identify where that paper contains a longer proof of a shared
compression calculation. The new continuous-domain reductions, posterior
arguments, teacher-tube construction, and approximate-readout duality are
proved below.

\section{Model and main results}\label{sec:model}
Fix $d\geq1$, $L\geq2$, and a nonempty compact set
$X\subseteq\{x\in\R^d:\|x\|_2\leq1\}$. Write $\bar x=(1,x)$.
A network of width $N$ and depth $L$ computes
\[
 h^0(x)=\bar x,\qquad h^\ell(x)=\tanh(W^\ell h^{\ell-1}(x))
 \quad(1\leq\ell<L),\qquad z_w(x)=v^Th^{L-1}(x).
\]
All hidden layers have width $N$, there are no separate internal biases,
and the readout is linear. Let $f:X\to\R$ be computed by a teacher of
the same depth, width $C\geq1$, and weights of absolute value at most
$B\geq1$. Width and weight magnitude are separate parameters. The theorem
does not claim that width alone controls arbitrary unbounded-weight teachers.

For a raw parameter prior $\mu$, inputs $x_i$ sampled independently from
any Borel probability distribution $D$ on $X$, and exact labels $f(x_i)$,
define the posterior by
\begin{equation}\label{eq:posterior}
 \Pi(dw\mid S_n)=
 \frac{\prod_{i=1}^n\mathbf1\{|z_w(x_i)-f(x_i)|<\rho\}\,\mu(dw)}
 {\int\prod_{i=1}^n\mathbf1\{|z_u(x_i)-f(x_i)|<\rho\}\,\mu(du)},
 \qquad \rho>0.
\end{equation}
There is no prior conditioning on unobserved inputs. The denominator is
positive because a neighborhood of the teacher clone fits uniformly.
For $\tau>\rho$, put
\[
 R_\tau(\Pi)=\E_{w\sim\Pi}D\{|z_w-f|>\tau\},\qquad
 q_\tau(x)=\Pi\{|z_w(x)-f(x)|\leq\tau\mid S_n\}.
\]
Tonelli and Markov give the posterior-confidence implication
\begin{equation}\label{eq:confidence}
 D\{q_\tau\leq2/3\}\leq3R_\tau(\Pi).
\end{equation}
The result controls the whole posterior, including mixtures, rather than
only its mean output or a selected parameter minimizer.

Our coordinate law $\nu$ has an even variance-one density
\begin{equation}\label{eq:law}
 q(t)=Z^{-1}e^{-V(t)},\qquad V\in C^2(\R),\qquad
 0<m\leq V''(t)\leq M<\infty.
\end{equation}
The law is fixed as width varies. Standard Gaussian coordinates are the
case $m=M=1$. A non-Gaussian example is the variance-standardized density
proportional to $\exp\{-t^2/2-b\log\cosh t\}$ for any fixed $b>0$.
All scalar coordinates below are independent copies of $\xi\sim\nu$.
The strong-prior scales are
\begin{equation}\label{eq:strong-scales}
 W^1_{ij}=\sigma\sqrt{N/d}\,\xi^1_{ij},\qquad
 W^\ell_{ij}=\sigma\xi^\ell_{ij}\ (2\leq\ell<L),\qquad
 v_j=\sigma N^{-1/2}\xi^L_j.
\end{equation}
The finite-scale family sets $\sigma=N^{-\gamma/2}$:
\begin{equation}\label{eq:finite-scales}
 W^1_{ij}=N^{(1-\gamma)/2}\xi^1_{ij}/\sqrt d,\quad
 W^\ell_{ij}=N^{-\gamma/2}\xi^\ell_{ij},\quad
 v_j=N^{-(\gamma+1)/2}\xi^L_j.
\end{equation}
At $\gamma=1$, the three variances are $1/d$, $1/N$, and $1/N^2$.
For compact notation write
\[
 u=\frac1{\min\{1,\tau-\rho\}},\qquad
 v_\rho=\log\frac1{\min\{1,\rho\}}.
\]

\begin{theorem}[Strong-prior regression]\label{thm:strong}
For every fixed $L$ and coordinate law~\eqref{eq:law}, there is a
polynomial $p_L$ such that the following holds. For every teacher as above,
every $N\geq C$, every $0<\rho<\tau$, and every $\varepsilon\in(0,1)$,
there is $\sigma_0>0$ such that, for all $0<\sigma<\sigma_0$, all
distributions $D$ on $X$, all $\delta\in(0,1)$, and
\[
 n\geq\frac6\varepsilon
 \left[p_L(d,C,B,u)+\log\frac1\delta\right],
\]
the posterior~\eqref{eq:posterior} satisfies, with probability at least
$1-\delta$ over the sample,
\[
 R_\tau(\Pi)\leq\varepsilon/3,\qquad
 D\{q_\tau\leq2/3\}\leq\varepsilon.
\]
The polynomial is independent of learner width. The scale threshold may
depend on that finite width, the teacher, $\rho$, and $\varepsilon$;
it does not depend on $D,n,\delta$.
\end{theorem}

\begin{theorem}[Finite-scale regression, including mean field]
\label{thm:finite}
Fix $L$, a coordinate law~\eqref{eq:law}, and $\gamma\geq1$.
There is a polynomial $p_{L,\gamma}$ such that, for every teacher as above,
every $N\geq C$, every $D$ on $X$, all $0<\rho<\tau$, and
$\varepsilon,\delta\in(0,1)$, the condition
\begin{equation}\label{eq:finite-samples}
 n\geq\frac6\varepsilon\left[
 p_{L,\gamma}\!\left(d,C,B,u,v_\rho,\log\frac1\varepsilon\right)
 +\log\frac6{\varepsilon\delta}\right]
\end{equation}
implies the same risk and posterior-confidence bounds with sample
probability at least $1-\delta$.
The polynomial is independent of $N$. Its degree and constants may depend
on $L,\gamma$ and the coordinate law; no uniformity as
$\gamma\downarrow1$ is claimed for the $\gamma>1$ proof.
The endpoint has a separate proof.
\end{theorem}

These theorems are finite-width statistical statements. They assert no
efficient method for sampling the posterior, no polynomial dependence on
depth, and no theorem for $\gamma<1$. The polynomial degree can itself grow
quickly with depth; we do not assert a bound with merely an exponential
depth prefactor. A different fixed confidence level
below one follows by replacing the factor $3$ in~\eqref{eq:confidence}.

\begin{remark}[Any bounded coordinate domain]\label{rem:cube}
The same conclusions hold for every compact $X\subseteq[-1,1]^d$.
All stability and rounding estimates below already use only coordinate
bounds. In the input-net estimates replace the unit-ball radius by
$\sqrt d$, so a mesh $r$ needs at most $(1+2\sqrt d/r)^d$ points.
Its logarithm still has a polynomial majorant. This includes finite input
sets as well as balls and spheres.
\end{remark}

\section{Two posterior estimates}\label{sec:posterior}
Let
\begin{equation}\label{eq:tube}
 U_\rho=\{w:\sup_{x\in X}|z_w(x)-f(x)|<\rho\}.
\end{equation}
This is an open parameter event, since $X$ is compact. It lower-bounds
the posterior denominator for every sample.

\begin{lemma}[Learning from a finite approximate cover]\label{lem:posterior}
Suppose $\mu(U_\rho)>0$, a measurable event $G$ obeys
$\mu(G^c)/\mu(U_\rho)\leq b$, and every output from $G$ is within
$a$ uniformly on $X$ of a member of a deterministic finite family
$\mathcal H$. If $2a<\tau-\rho$ and
\[
 n\geq t^{-1}[\log|\mathcal H|+\log(1/\delta)],
\]
then $R_\tau(\Pi)\leq t+b$ with sample probability at least $1-\delta$.
\end{lemma}
\begin{proof}
For each $h\in\mathcal H$, set
$A_h=\{x:|h(x)-f(x)|>\rho+a\}$. If $D(A_h)>t$, its probability
of containing no training input is at most $e^{-nt}$. A union bound
excludes all such surviving $h$ with the asserted probability.
Every fitted $w\in G$ has an approximant $h$ whose set $A_h$ contains
no training input. Moreover $|z_w-f|>\tau$ implies $x\in A_h$,
by $\tau-a>\rho+a$. Thus its test risk is at most $t$.
For every sample, $\Pi(G^c\mid S_n)\leq\mu(G^c)/\mu(U_\rho)\leq b$.
Integrating establishes the conclusion.
\end{proof}

\begin{lemma}[Learning from a uniform tube]\label{lem:mass-posterior}
If $\mu(U_\rho)\geq e^{-A}$, then
$n\geq t^{-1}[A+\log(1/(b\delta))]$ implies
$R_\tau(\Pi)\leq t+b$ with sample probability at least $1-\delta$.
\end{lemma}
\begin{proof}
A fixed parameter with test risk greater than $t$ fits all stricter
training constraints with probability at most $e^{-nt}$. Tonelli gives
expected raw mass of such fitted parameters at most $e^{-nt}$.
Divide by the deterministic denominator lower bound $e^{-A}$ and apply
Markov to bound their posterior mass by $b$ with probability $1-\delta$.
The remaining parameters each have risk at most $t$.
\end{proof}

All covers constructed below are finite collections of Borel functions.
The relevant exceptional events are measurable: for each such function,
the uniform distance from a network output is a Borel function of the
parameters, and the constructions alternatively supply explicit measurable
norm and covariance-certificate events. No measurable choice of approximant
or of teacher-network parameters is needed.

\section{Cloning and the mass of a teacher tube}\label{sec:mass}
Define the weighted energy
\begin{equation}\label{eq:energy}
 \kappa_N(w)=\frac dN\|W^1\|_F^2+
       \sum_{\ell=2}^{L-1}\|W^\ell\|_F^2+N\|v\|_2^2.
\end{equation}
Partition each learner layer into $C$ groups with balanced sizes
$k_j\in\{\lfloor N/C\rfloor,\lceil N/C\rceil\}$. Copy the
teacher neuron in every position of its group, divide an internal teacher
edge by its source-group size, and divide each readout edge by its group
size. The clone computes $f$ exactly for every real input. Since
$\max k_i/\min k_i\leq2$ and $N/\min k_i\leq2C$, its energy is at most
\begin{equation}\label{eq:teacher-energy}
 E_*:=d(d+1)B^2+2(L-1)C^2B^2.
\end{equation}
In particular, the first-layer contribution is at most $d(d+1)B^2$,
each internal contribution at most $2C^2B^2$, and the readout contribution
at most $2C^2B^2$. These estimates are independent of $N$.

We use three consequences of~\eqref{eq:law}. A product of $\nu$
coordinates conditioned on a centrally symmetric convex set is centered
and has
\begin{equation}\label{eq:subgaussian}
 \E e^{t\langle u,Z\rangle}\leq e^{t^2\|u\|^2/(2m)}.
\end{equation}
Indeed convex conditioning preserves $m$-strong log-concavity. The
Brascamp--Lieb variance inequality after every linear tilt bounds the
Hessian of the log moment generating function by $m^{-1}I$; symmetry
sets its derivative at zero to zero. Smooth convex barriers justify
conditioning on sets. For the unconditioned product, Caffarelli's
contraction theorem and Gaussian isoperimetry give
\begin{equation}\label{eq:enlargement}
 \mu(A_r)\geq\Phi(\Phi^{-1}(\mu(A))+\sqrt m\,r).
\end{equation}
Here $A_r$ is Euclidean enlargement. Finally any symmetric event $E$ and
translation $a$ satisfy
\begin{equation}\label{eq:translation}
 \mu(E+a)\geq\mu(E)e^{-M\|a\|^2/2}.
\end{equation}
To prove the last inequality, translate the conditional law on $E$.
Its relative entropy against $\mu$ is
$-\log\mu(E)+\E[\sum_i(V(Z_i+a_i)-V(Z_i))\mid E]$.
Taylor's theorem bounds the added expectation by $M\|a\|^2/2$:
the linear score term vanishes by oddness of $V'$ and symmetry.
A law supported on $E+a$ has relative entropy at least $-\log\mu(E+a)$.
References for the two concentration theorems are
\cite{brascamp-lieb,caffarelli,ledoux}.

\begin{proposition}[A uniform teacher tube at every width]\label{prop:tube}
For $N\geq C$ and $\gamma\geq1$,
\begin{equation}\label{eq:tube-mass}
 \mu_{N,\gamma}(U_\rho)\geq e^{-A_L(d,C,B,\rho)N^\gamma},
\end{equation}
where $A_L$ has a polynomial majorant in $d,C,B,v_\rho$ at fixed
depth and coordinate law.
\end{proposition}
\begin{proof}
Let $P$ project onto the group-constant subspace of the clone, with basis
$u_j=k_j^{-1/2}\mathbf1_{\text{group }j}$. Put
$t=N^{(1-\gamma)/2}\leq1$. In standardized coordinates impose the
centered cylinder consisting of $|\xi|\leq\eta$ for first-layer and
readout entries, and $|\langle Z_{i,*},u_j\rangle|\leq\eta$ for every
internal row and group. These scalar constraints concern independent
variables, because groups are disjoint.

Every normalized group sum is even, log-concave, and has variance at most
$1/m$. Symmetric unimodality and Chebyshev yield, for $0<\eta\leq1$,
\[
 \Prb(|Y|\leq\eta)\geq c_m\eta,
 \qquad c_m=\tfrac38\min\{1,\sqrt m\}.
\]
For example take the interval of radius
$2\max\{1,m^{-1/2}\}$, which has mass at least $3/4$, and compare
average densities on concentric intervals. Hence the cylinder has mass
at least $(c_m\eta)^{N[d+2+(L-2)C]}$.

Conditional on the cylinder,~\eqref{eq:subgaussian} and two $1/4$ sphere
nets of size at most $9^N$ give
\[
 \Prb(\|Z\|_{\rm op}>2q\sqrt N\mid\text{cylinder})
 \leq2e^{(2\log9-mq^2/2)N}.
\]
Set $b_0=2\sqrt{(2/m)[2\log9+\log(4L)]}$. Intersect the cylinder
with $\|Z\|_{\rm op}\leq b_0\sqrt N$ for every internal layer.
The resulting symmetric event $E$ has at least half the cylinder's mass.

Translate by the standardized clone. The internal perturbation is
$tZ/\sqrt N=tZP/\sqrt N+tZ(I-P)/\sqrt N$. Its second summand
annihilates the teacher activations, and its first has norm at most
$\eta\sqrt C$. If $e_\ell=N^{-1/2}\|h^\ell-h_*^\ell\|_2$,
boundedness and contraction of tanh give, uniformly on $X$,
\[
 e_1\leq\eta(d+1),\qquad
 e_\ell\leq(\sqrt2CB+b_0)e_{\ell-1}+\eta\sqrt C.
\]
Put $A_0=2+3CB+b_0$. Then
$e_{L-1}\leq\eta(d+L+1)A_0^{L-2}$. The clone has
$\sqrt N\|v_*\|\leq\sqrt2CB$; the perturbation readout has
$\ell^1$ norm at most $\eta$. Its output error is therefore at most
$\eta P_0$, where
\[
 P_0=1+2CB(d+L+1)A_0^{L-2},\qquad
 \eta=\min\{1/2,\rho/(2P_0)\}.
\]
The translated event lies in $U_\rho$. Its standardized translation has
squared norm $N^\gamma\kappa_N(w_*)\leq E_*N^\gamma$.
Equation~\eqref{eq:translation} and $N\leq N^\gamma$ prove the claim
with, for instance,
\begin{equation}\label{eq:A}
 A_L=1+\tfrac M2E_*+[d+2+(L-2)C]\log\frac1{c_m\eta}+\log2.
\end{equation}
All stability estimates used only bounded input coordinates, so they
hold simultaneously on the continuous domain.
\end{proof}

\section{Uniform covers for bounded energy}\label{sec:energy-cover}
The following finite-input calculation is the deterministic compression
ingredient shared with~\cite{nongaus}. Only the logarithm of the number
of input points enters it.

\begin{lemma}[Finite-input neuron compression]\label{lem:compression}
Let $T\subset[-1,1]^d$ be finite of cardinality $s$, $K\geq1$, and
$0<a\leq1$. Every network with $\kappa_N(w)\leq K$ has a
width-$J$ approximant with energy at most $4K$ and error at most $a/4$
on $T$, where $J$ has a polynomial majorant in
$d,K,1/a,\log(2s)$ at fixed $L$. One may take
\[
 \lambda=\frac{a2^{-L}K^{-(L-1)/2}}{4(L-1)},\quad
 D_T=\max\{1,\lceil\log_2s\rceil\},\quad
 J=\max\left\{\lceil4K\rceil,
 \left\lceil256\lambda^{-4}[D_T+\log(L/\lambda)]\right\rceil\right\}.
\]
All weights of the approximant have magnitude at most $J$.
\end{lemma}
\begin{proof}
Here are the estimates underlying the neuron-sampling construction;
the full reverse-order telescoping calculation is also given in the
bounded-energy compression appendix of~\cite{nongaus}.
In each internal or readout row of Euclidean norm $r_i$, delete entries
exceeding $2r_i/(\lambda\sqrt N)$. Their total absolute contribution
against a bounded hidden vector is at most
$(\lambda/2)\sqrt N r_i$. At every hidden layer independently sample
$J$ uniform neuron indices with replacement. Retain the sampled first-layer
rows; multiply each retained edge out of a sampled neuron by $N/J$.

For each retained row and input, Hoeffding's inequality bounds the
probability that the estimated sum differs from the truncated sum by
more than $(\lambda/2)\sqrt N r_i$ by
$2e^{-\lambda^4J/32}$. Union bounding over inputs, rows, and layers
costs at most $2sLJ e^{-\lambda^4J/32}\leq1/100$.
The expected weighted energy of the sampled network is at most $K$;
independence of adjacent layers' index lists suffices for the internal
matrix calculation. Thus its energy is at most $4K$ with probability
at least $3/4$. The sum of the squared norms of retained rows, with
the $N$ normalization, is at most $8JK$ with probability at least $7/8$,
again by Markov. These three events have a nonempty intersection.

Fix a realization in that intersection and expose sampled layers in
reverse order. The downstream matrices at each step already equal their
final sampled matrices, whose norms are at most $2\sqrt K$, and the
sampled readout has norm at most $2\sqrt{K/J}$. Propagating the row
errors shows that each truncation and each sampling stage changes the
output by at most $a/[8(L-1)]$. The $2(L-1)$ stages total at most
$a/4$. This is precisely why the choice of $\lambda$ contains
$2^{-L}K^{-(L-1)/2}$. The energy bound gives entrywise bounds
$\sqrt{4KJ/d}$, $2\sqrt K$, and $2\sqrt{K/J}$ at the three layer
types, each at most $J$ since $J\geq4K$.
\end{proof}

\begin{proposition}[A continuous-domain energy cover]\label{prop:energy-cover}
For fixed $L$, $K\geq1$, and $0<a\leq1$, there is a deterministic finite
family $\mathcal H_{K,a}$ of Borel functions on $X$, common to all
learner widths, such that
\[
 \inf_{h\in\mathcal H_{K,a}}\sup_X|z_w-h|\leq a
 \quad\text{whenever }\kappa_N(w)\leq K,
\]
and $\log|\mathcal H_{K,a}|\leq H_L(d,K,1/a)$ for a polynomially
bounded function at fixed depth.
\end{proposition}
\begin{proof}
Multiplying the layer norm bounds from~\eqref{eq:energy} gives
\begin{equation}\label{eq:energy-lipschitz}
 \Lip(z_w)\leq K^{L/2}/\sqrt d.
\end{equation}
Choose an input net $T\subset X$ of radius
$r=\min\{1,a\sqrt d/(8K^{L/2})\}$ and size
$s\leq(1+2/r)^d$. Fix a nearest-net-point map $t:X\to T$, resolving
ties by a fixed enumeration. It is Borel. Apply
Lemma~\ref{lem:compression} on $T$.

For any width-$J$ network with weights bounded by $J$, rounding every
weight to a grid of mesh
\[
 \Delta\leq\frac{a}{4(d+L+1)(2J^2)^L}
\]
changes its output by at most $a/4$ whenever the input coordinates are
bounded by one. To check this, a first-layer coordinate error is at most
$(d+1)\Delta$; each next coordinate error is at most $2J^2$ times
the preceding error plus $J\Delta$, including the readout. Induction
gives the displayed bound.
There are $P=J(d+1)+(L-2)J^2+J$ coordinates, so at most
$(1+2(J+1)/\Delta)^P$ rounded networks.

Use their functions $x\mapsto z_{\rm rounded}(t(x))$ as cover centers.
Compression, rounding, and~\eqref{eq:energy-lipschitz} give total error
at most $a/4+a/4+a/8<a$. The centers need not themselves be continuous
or neural networks; they are used only in the statistical union bound.
Since $\log s=O_L(d\log(2+K+1/a))$, the logarithm of their number has
a fixed-depth polynomial majorant and is independent of $N$.
\end{proof}

\section{The strong-prior and super-mean-field proofs}\label{sec:strong-above}
\begin{proof}[Proof of Theorem~\ref{thm:strong}]
Fix $N$ and the exact clone $w_*$. The tube $U_\rho$ is an open
neighborhood of $w_*$. Evenness and~\eqref{eq:law} give two-sided
quadratic density envelopes
$c e^{-Mt^2/2}\leq q(t)\leq C_0e^{-mt^2/2}$.
In the finite-dimensional parameter space, elementary Gaussian-integral
bounds therefore yield
\[
 \liminf_{\sigma\downarrow0}\sigma^2\log\mu_\sigma(U_\rho)
 \geq-\tfrac M2E_*,\qquad
 \limsup_{\sigma\downarrow0}\sigma^2\log
       \mu_\sigma\{\kappa_N>K\}\leq-\tfrac m2K.
\]
For the lower bound integrate a sufficiently small open neighborhood of
$w_*$ whose energy is at most $E_*+\zeta$, then let $\zeta\downarrow0$.
For the upper bound integrate the quadratic envelope outside the ellipsoid;
fixed dimension contributes only subexponential factors.
Take $K=(M/m)E_*+2$. The exceptional-to-tube mass ratio tends to zero.
Choose $\sigma_0$ making it at most $\varepsilon/6$ for all smaller
scales. Apply Proposition~\ref{prop:energy-cover} with
$a=\min\{1/2,(\tau-\rho)/4\}$ and then
Lemma~\ref{lem:posterior} with $t=b=\varepsilon/6$.
Equation~\eqref{eq:confidence} finishes the proof.
\end{proof}

The same proof uses only independent coordinate densities with common
two-sided envelopes $c_j e^{-c_+t^2}\leq q_j(t)\leq C_j e^{-c_-t^2}$,
where $0<c_-\leq c_+<\infty$ do not depend on width. Symmetry and
log-concavity are unnecessary in this strong-prior extension. Replace
$K$ by $(c_+/c_-)E_*+2$; prefactors affect only the fixed-width scale
threshold.

\begin{lemma}[Spectral truncation at $\gamma>1$]\label{lem:spectral}
Fix $K\geq1$ and $0<a<1$, and put
$b=a/[2L(\sqrt K+1)^{L-2}]$. If
\begin{equation}\label{eq:above-threshold}
 N^{\gamma-1}\geq8(d+1)/m,\qquad
 N^{\gamma-1}\geq256/(mb^2),
\end{equation}
then, except with probability $Le^{-mKN^\gamma/8}$, the network output
is uniformly within $a/2$ on $X$ of a network of energy at most $LK$.
\end{lemma}
\begin{proof}
Equation~\eqref{eq:subgaussian} implies
$\E e^{m\xi^2/4}\leq\sqrt2$ by integrating against a one-dimensional
Gaussian. Hence, for $D_0$ independent coordinates,
\begin{equation}\label{eq:squares}
 \Prb\left(\sum_{j=1}^{D_0}\xi_j^2>u\right)
 \leq2^{D_0/2}e^{-mu/4}\leq e^{-mu/8}
 \quad\text{if }mu\geq8D_0.
\end{equation}
This controls the first-layer weighted norm and the readout weighted norm
by $K$ outside two events of probability $e^{-mKN^\gamma/8}$.

For an internal matrix $A=\sigma Z$, the net bound gives
$\Prb(\|A\|_{\rm op}\leq8\sigma\sqrt{N/m})\geq1/2$.
If $b\geq16\sigma\sqrt{N/m}$, the operator-norm ball of radius $b/2$
has mass at least $1/2$. Its Frobenius distance from $A$ is
$D(A)^2=\sum_j(s_j(A)-b/2)_+^2$. Equation~\eqref{eq:enlargement}
implies $\Prb(D(A)>r)\leq e^{-mr^2/(2\sigma^2)}$.
The singular part $A_{>b}$ has $\|A_{>b}\|_F\leq2D(A)$, whence
\[
 \Prb(\|A_{>b}\|_F^2>K)\leq e^{-mK/(8\sigma^2)}.
\]
At $\sigma=N^{-\gamma/2}$ the second condition
in~\eqref{eq:above-threshold} permits this estimate.

Discard all internal singular values at most $b$. The retained network
has energy at most $LK$. If $e_\ell$ denotes its Euclidean hidden error,
$e_1=0$ and $e_\ell\leq\sqrt K e_{\ell-1}+b\sqrt N$.
The readout norm is at most $\sqrt{K/N}$, so the uniform output error
is at most $b\sum_{j=1}^{L-2}K^{j/2}\leq a/2$.
The $L=2$ sum is empty. A union bound proves the assertion.
\end{proof}

\begin{proof}[Proof of Theorem~\ref{thm:finite} when $\gamma>1$]
Let $A=A_L(d,C,B,\rho)$, $a=\min\{1/2,(\tau-\rho)/4\}$, and
\[
 K=\max\{1,(8/m)[A+\log(6L/\varepsilon)]\}.
\]
Let $T$ be the ceiling of the maximum of $C$ and the two thresholds
in~\eqref{eq:above-threshold}. For $N\geq T$,
Lemma~\ref{lem:spectral} and a cover of the retained networks to accuracy
$a/2$ give a common finite output cover of accuracy $a$.
By~\eqref{eq:tube-mass}, its exceptional-to-tube mass ratio is at most
$Le^{-N^\gamma(mK/8-A)}\leq\varepsilon/6$.
Apply Lemma~\ref{lem:posterior}.
For $C\leq N<T$, use Lemma~\ref{lem:mass-posterior} with tube cost
at most $AT^\gamma$. At fixed $L,\gamma>1$, the cover entropy and
$AT^\gamma$ have polynomial majorants in the arguments
of~\eqref{eq:finite-samples}. This covers every $N\geq C$.
\end{proof}

\section{The mean-field endpoint}\label{sec:meanfield}
At $\gamma=1$ the typical internal matrices have nonvanishing operator
norm. Discarding their bulk does not produce a good approximation.
Instead, finite covariance tests describe most prior outputs, and a
bounded dual representation turns covariance approximation into output
approximation. The argument below retains the actual first-layer law.
Auxiliary Gaussian rows are used to construct finite descriptions; the
learner's posterior is not replaced by a Gaussian process.

\subsection{A bounded dual for an arbitrary readout}
\begin{lemma}[Approximate-readout duality]\label{lem:readout}
Let $H\in\R^{s\times N}$, $Q=HH^T/N$, $z=Hv$, and
$N\|v\|_2^2\leq K$. For every $t>0$ there exists $\beta\in\R^s$
such that
\[
 \|Q\beta-z\|_\infty\leq t,\qquad
 \|\beta\|_1\leq\frac K{4t}.
\]
\end{lemma}
\begin{proof}
Minimize $N\|u\|_2^2/2$ subject to $\|Hu-z\|_\infty\leq t$.
The point $v$ is strictly feasible. Convex optimality gives a signed
difference $\beta$ of the lower and upper multipliers with
$Nu=H^T\beta$. Complementarity gives
\[
 t\|\beta\|_1=\beta^T(z-Hu)
 =Nv^Tu-N\|u\|^2.
\]
The two multipliers for the same coordinate cannot both be positive,
because $t>0$. Put $r=\sqrt N\|u\|$. The right side is at most
$\sqrt K r-r^2\leq K/4$. Also $Hu=Q\beta$, proving both claims.
Neither a sign margin nor invertibility of $Q$ is required.
\end{proof}

\subsection{Finite covariance certificates}
For a law $\mu$ on $[-1,1]^s$, a symmetric matrix $\Gamma$, and an
empirical covariance $\widehat Q$, write
\[
 \mathcal J_G(\widehat Q;\mu)=
 \max_{\|\Gamma\|_{1,\rm ent}\leq G}
 \left\{\langle\Gamma,\widehat Q\rangle
       -\log\E_\mu e^{h^T\Gamma h}\right\},\qquad
 \|\Gamma\|_{1,\rm ent}=\sum_{i,j}|\Gamma_{ij}|.
\]
The norm dual to this entrywise $\ell^1$ norm is
$\|A\|_{\max}=\max_{i,j}|A_{ij}|$.

\begin{lemma}[Distribution-free certificate]\label{lem:certificate}
For an empirical second moment of $N$ independent samples from $\mu$,
\begin{equation}\label{eq:certificate}
 \Prb\{\mathcal J_G(\widehat Q;\mu)>K\}
 \leq\exp\{q\log(2s^2+1)-N(K-1)\},
 \qquad q=\lceil G^2/2\rceil.
\end{equation}
If $\mathcal J_G\leq K$, some quadratic exponential tilt of $\mu$,
with parameter norm at most $G$, has covariance within $K/G$ of
$\widehat Q$ in maximum norm.
\end{lemma}
\begin{proof}
Flatten the $s^2$ bounded observables $F_{ij}=h_ih_j$ into a vector
$F\in[-1,1]^D$, and write
$\Lambda(\lambda)=\log\E e^{\lambda\cdot F}$.
Every $\ell^1$-ball vector $\lambda$ is the mean of a random vector
in $\{0,\pm Ge_1,\ldots,\pm Ge_D\}$. Average $q$ independent such
vectors to get a finite set $\mathcal V$ of size at most $(2D+1)^q$.
Hoeffding's exponential-moment bound and Jensen imply
$\E\Lambda(V)\leq\Lambda(\lambda)+G^2/(2q)$.
Thus for every moment vector $a$,
\[
 \max_{\|\lambda\|_1\leq G}\{\lambda\cdot a-\Lambda(\lambda)\}
 \leq1+\max_{v\in\mathcal V}\{v\cdot a-\Lambda(v)\}.
\]
For empirical $a$, independence gives
$\E e^{N[v\cdot a-\Lambda(v)]}=1$. Markov and the finite union
bound prove~\eqref{eq:certificate}.

For the approximation assertion, let $\lambda_*$ maximize the dual and
let $\mu_*$ be the corresponding exponential tilt, with moments $a_*$.
The first-order condition for maximization over the $\ell^1$ ball gives
$\lambda_*\cdot(a-a_*)=G\|a-a_*\|_\infty$. Therefore
\[
 \mathcal J_G=\KL(\mu_*\Vert\mu)+G\|a-a_*\|_\infty.
\]
Relative entropy is nonnegative. Symmetrizing the parameter matrix does
not change its quadratic form or increase its norm.
\end{proof}

For $Q\succeq0$, let $\mathsf F_\Gamma(Q)$ be the covariance of
$h=\tanh Z$, $Z\sim\mathcal N(0,Q)$, tilted by $e^{h^T\Gamma h}$.
Two bounds control this Gaussian auxiliary family.

\begin{lemma}[Relative comparison and covariance stability]
\label{lem:comparisons}
Let $\|a_j\|_\infty\leq N^{-1/2}$, let $\xi_j\sim\nu$ and
$g_j\sim\mathcal N(0,1)$ independently, and put
$Z_\nu=\sum_j a_j\xi_j$, $Z_G=\sum_j a_jg_j$.
For $\|\Gamma\|_{1,\rm ent}\leq G$ and $N\geq4G^2$,
\begin{equation}\label{eq:lindeberg}
 \left|\log\E e^{\tanh(Z_\nu)^T\Gamma\tanh(Z_\nu)}
      -\log\E e^{\tanh(Z_G)^T\Gamma\tanh(Z_G)}\right|
 \leq\frac{c_\nu P_3(G)}{\sqrt N},
\end{equation}
where $P_3(G)=24G+36G^2+8G^3$ and
$c_\nu=(e/6)(\E|\xi|^3e^{|\xi|}+\E|g|^3e^{|g|})<\infty$.
Also, for all positive semidefinite $Q,Q'$,
\begin{equation}\label{eq:covariance-stability}
 \|\mathsf F_\Gamma(Q)-\mathsf F_\Gamma(Q')\|_{\max}
 \leq8(1+G)^2\|Q-Q'\|_{\max}.
\end{equation}
Neither bound depends on the number $s$ of field coordinates.
\end{lemma}
\begin{proof}
Put $F(z)=\tanh(z)^T\Gamma\tanh(z)$. In a direction
$\|a\|_\infty\leq N^{-1/2}$ its first three directional derivatives
are bounded by $2G/\sqrt N$, $6G/N$, and $24G/N^{3/2}$.
Consequently $|D_a^3e^F|\leq P_3(G)N^{-3/2}e^F$.
At a one-coordinate replacement, fix the residual $z$ and define
$f_z(t)=e^{F(z+at)}$. Since $2G/\sqrt N\leq1$,
\[
 e^{-|t|}f_z(0)\leq f_z(t)\leq e^{|t|}f_z(0).
\]
Taylor expansion through order two and matched mean and variance imply
that the two averaged values differ by at most
$P_3(G)N^{-3/2}(\E|\xi|^3e^{|\xi|}+\E|g|^3e^{|g|})f_z(0)/6$.
Each average is at least $e^{-1}f_z(0)$ by Jensen and variance one.
These bounds persist after averaging over the residual. The inequality
$|\log u-\log v|\leq|u-v|/\min(u,v)$ and $N$ replacements give
\eqref{eq:lindeberg}. This relative calculation avoids an exponential
loss in the tilt budget.

For~\eqref{eq:covariance-stability}, interpolate covariance linearly
and differentiate the Gaussian tilted expectation of
$H(z)=\tanh(z_i)\tanh(z_j)$, using the quotient rule. If
$M_k$ denotes the supremum of the sum of absolute derivatives of order
$k$, then $M_1(F)\leq2G$, $M_2(F)\leq6G$,
$M_1(H)\leq2$, and $M_2(H)\leq6$. Gaussian differentiation bounds
the derivative by
\[
 \tfrac12\|Q-Q'\|_{\max}
       [6+8G+2(6G+4G^2)]
 \leq8(1+G)^2\|Q-Q'\|_{\max}.
\]
All terms are expectations or covariances under the normalized tilted
law, so no reciprocal partition-function estimate is used. Integrate
the derivative. Adding $\eta I$ and taking $\eta\downarrow0$
justifies singular covariance matrices. Expanded calculations for both
estimates appear in~\cite{nongaus}.
\end{proof}

\subsection{Most outputs have a common finite cover}
\begin{proposition}[Mean-field output cover]\label{prop:mf-cover}
There are fixed-depth polynomially bounded functions
$T_L(d,K,1/a)$ and $H_L^{\rm mf}(d,K,1/a)$ with the following
property. For $0<a<1$ and $K$ above a fixed-law constant at least $4$,
there is a finite deterministic family $\mathcal C_{K,a}$, independent
of learner width, such that for every $N\geq T_L(d,K,1/a)$,
\[
 \log|\mathcal C_{K,a}|\leq H_L^{\rm mf}(d,K,1/a),
 \qquad
 \mu_{N,1}\!\left\{\inf_{h\in\mathcal C_{K,a}}\sup_X|z_w-h|>a\right\}
 \leq2Le^{-cNK},
\]
where $c=\min\{1/4,m/16\}$.
\end{proposition}
\begin{proof}
We give the construction, including its width and net dependencies.

\emph{Lipschitz event.}
Require
\[
 \|W^1\|_F^2\leq NK,\qquad
 \|W^\ell\|_{\rm op}\leq\sqrt K\ (2\leq\ell<L),\qquad
 N\|v\|^2\leq K.
\]
On this event $\Lip(z_w)\leq K^{L/2}$. The square moment
bound~\eqref{eq:squares} gives first-layer and readout failure at most
$e^{-mNK/8}$ for sufficiently large fixed-law $K$. For the first layer
the exponent before simplification is
$-mdNK/4+N(d+1)\log2/2$.
The sphere-net estimate gives internal failure at most
$2e^{2N\log9-mNK/8}\leq e^{-mNK/16}$ after enlarging the same
constant baseline. Thus total failure is at most $Le^{-mNK/16}$.

Choose a net $S\subset X$ of radius $a/(4K^{L/2})$ and cardinality
\begin{equation}\label{eq:mf-net}
 s\leq(1+8K^{L/2}/a)^d.
\end{equation}
Fix a Borel nearest-net-point map $t:X\to S$. The actual output changes
by at most $a/4$ between $x$ and $t(x)$ on the Lipschitz event.

\emph{Proxy covariance tower.}
Write $r=L-1$, $\bar K=K+1$, and
\[
 G_r=8rK\bar K/a^2,\qquad
 G_\ell=8G_{\ell+1}(1+G_{\ell+1})^2\quad(\ell<r),\qquad
 b_\ell=8(1+G_\ell)^2.
\]
Let $\ell_0=\log3+2d\log(1+8K^{L/2}/a)$, so that
$\log(2s^2+1)\leq\ell_0$ uniformly over compact $X$ in the unit ball. For
$X\subseteq[-1,1]^d$ as in Remark~\ref{rem:cube}, replace $8K^{L/2}/a$ by
$8\sqrt d\,K^{L/2}/a$ here and in~\eqref{eq:mf-net}. Take
\begin{equation}\label{eq:T}
 T_L=\left\lceil\max\left\{1,
 \frac{2\lceil G_1^2/2\rceil\ell_0}{K-1},
 4G_1^2,c_\nu^2P_3(G_1)^2\right\}\right\rceil.
\end{equation}
These budgets and this threshold have polynomial majorants in
$d,K,1/a$ at fixed $L$, using~\eqref{eq:mf-net}.

Let $Q_{\ell,N}$ be the empirical activation covariance indexed by $S$.
At the first layer, apply Lemma~\ref{lem:certificate} to its actual
non-Gaussian field law $\mu_0$. At later layers apply it conditionally
to the true row law. A union bound bounds failure by $r e^{-NK/4}$.
An internal preactivation is $N^{-1/2}\sum_j\xi_jh_j$, so
\eqref{eq:lindeberg} compares its log partition functions to the Gaussian
law of covariance $Q_{\ell-1,N}$ with error at most one under~\eqref{eq:T}.
The Gaussian certificate is therefore at most $\bar K$.
The lift in Lemma~\ref{lem:certificate} produces matrices $\Gamma_\ell$
with entrywise norm at most $G_\ell$. Define $Q_1^*$ by tilting the
actual first-layer law with $\Gamma_1$, and recursively set
$Q_\ell^*=\mathsf F_{\Gamma_\ell}(Q_{\ell-1}^*)$ for $\ell\geq2$.
Equation~\eqref{eq:covariance-stability} gives
\[
 \|Q_{r,N}-Q_r^*\|_{\max}
 \leq\sum_{\ell=1}^r\frac{\bar K}{G_\ell}
                     \prod_{j=\ell+1}^r b_j
 =\frac{r\bar K}{G_r}=\frac{a^2}{8K}.
\]

\emph{Readout and finite realization.}
On the readout norm event, Lemma~\ref{lem:readout} with tolerance $a/4$
gives $\|\beta\|_1\leq K/a$ and
$\|Q_{r,N}\beta-z_w|_S\|_\infty\leq a/4$.
Replacing the covariance by $Q_r^*$ adds error at most $a/8$.

Here is a finite network realizing the proxy to another $a/8$.
Put $b=8(1+G_1)^2$ and
\[
 \eta=\frac{a^2}{8Kr b^{r-1}},\qquad
 J=\left\lceil2\eta^{-2}\log(8rs^2)\right\rceil.
\]
Sample $J$ first-layer rows from the actual incoming law tilted by
$e^{h^T\Gamma_1h}$. At each later layer, sample independent Gaussian
rows divided by $\sqrt J$, tilted by
$e^{\tanh(Z)^T\Gamma_\ell\tanh(Z)}$ conditionally on the previous
sampled activations. Their reference Gaussian field has covariance equal
to the previous empirical covariance. Hoeffding and a union bound over
$rs^2$ covariance entries give simultaneous empirical error at most
$\eta$ with probability at least $3/4$.
Covariance stability then gives final covariance error at most
$r b^{r-1}\eta$. Set each readout coefficient to $(\beta\cdot h_i)/J$.
The output error against $Q_r^*\beta$ is at most
$(K/a)r b^{r-1}\eta=a/8$.

All sampled weights can simultaneously be bounded polynomially.
Each row's tilted density relative to its reference is at most $e^{2G_1}$.
The reference coordinates obey $\E e^{a_0\xi^2}\leq\sqrt2$ for
$a_0=\min\{m,1\}/4$. If $D_J=J(d+1)+(L-2)J^2$, choose
\[
 B_0=\max\left\{1,\frac K{aJ},
 \sqrt{\frac{2G_1+\log(4\sqrt2\max\{D_J,1\})}{a_0}}\right\}.
\]
A union bound gives probability at least $3/4$ that all incoming weights
are bounded by $B_0$. Internal weights are additionally divided by
$\sqrt J$, so the bound remains valid. The readout bound is deterministic.
Intersecting with the covariance event gives a realization with width and
weights polynomially bounded in $d,K,1/a$ at fixed depth.

\emph{Counting and extension off the net.}
For that common width and weight bound, put
\[
 P=J(d+1)+(L-2)J^2+J,\qquad
 \Lambda=(d+LJ)[1+J(B_0+1)]^{L-1}.
\]
Round parameters to mesh $a/(8\Lambda)$. A layer-by-layer
coordinate-error induction bounds the output change by $a/8$ on all
inputs with coordinates at most one. There are at most
$\exp\{P\log(3+16B_0\Lambda/a)\}$ rounded networks, with logarithm
polynomially bounded in the required parameters.

Include the functions $h(x)=z_{\rm rounded}(t(x))$ in the cover.
Its error to the actual output is at most
\[
 \underbrace{a/4}_{\text{readout dual}}
 +\underbrace{a/8}_{\text{proxy covariance}}
 +\underbrace{a/8}_{\text{finite realization}}
 +\underbrace{a/8}_{\text{rounding}}
 +\underbrace{a/4}_{\text{input net}}=7a/8<a.
\]
Only the actual output needs a Lipschitz bound; no circular choice of net
based on the approximant occurs. The norm and certificate failures total
at most $2Le^{-cNK}$. The grid of possible centers depends on the common
bounds and the fixed net, and hence is independent of $N$ and of the draw.
\end{proof}

\begin{proof}[Proof of Theorem~\ref{thm:finite} when $\gamma=1$]
Let $A=A_L(d,C,B,\rho)$ and $a=\min\{1/2,(\tau-\rho)/4\}$.
Choose $K$ above the baseline in Proposition~\ref{prop:mf-cover} and with
$cK\geq A+\log(12L/\varepsilon)$. At $N\geq T_L(d,K,1/a)$,
the exceptional-to-tube mass ratio is at most
$2Le^{-N(cK-A)}\leq\varepsilon/6$. Lemma~\ref{lem:posterior}
gives risk at most $\varepsilon/3$ at sample size
$(6/\varepsilon)[H_L^{\rm mf}(d,K,1/a)+\log(1/\delta)]$.
For $C\leq N<T_L$, Lemma~\ref{lem:mass-posterior} uses the tube
cost $AT_L$. Both bounds have the polynomial majorant
in~\eqref{eq:finite-samples}. Equation~\eqref{eq:confidence} completes
the argument at every width $N\geq C$.
\end{proof}


\section{Exact interpolation in an iterated strong-prior limit}
\label{sec:exact}
Positive training tolerance is necessary for the ordinary conditional
posterior~\eqref{eq:posterior} to be defined by a positive-probability
event. There is nonetheless a precise strong-prior exact-interpolation
conclusion. It follows from one sample event that controls all low-energy
interpolants, including all mixtures of them.

\begin{proposition}[Every low-energy interpolant generalizes]
\label{prop:exact}
Let $E_*$ be the clone energy bound~\eqref{eq:teacher-energy}, and fix
$\tau>0$. For fixed $L$, let $H$ be the logarithmic cardinality of the
all-width energy cover at energy $E_*$ and radius
$a=\min\{1,\tau/8\}$. It is polynomially bounded in
$d,C,B,1/\min\{1,\tau\}$. If
\[
 n\geq\frac3\varepsilon[H+\log(1/\delta)],
\]
then, with probability at least $1-\delta$, simultaneously for every
$N\geq C$ and every parameter satisfying
\[
 \kappa_N(w)\leq E_*,\qquad
 \max_i|z_w(x_i)-f(x_i)|\leq\tau/2,
\]
one has $D\{|z_w-f|>\tau\}\leq\varepsilon/3$.
In particular the same bound holds for every minimum-energy exact
interpolant and for every probability mixture of such interpolants.
All these mixtures satisfy~\eqref{eq:confidence}.
\end{proposition}
\begin{proof}
For a cover center $h$ set $A_h=\{|h-f|>3\tau/4\}$.
A covered parameter fitting within $\tau/2$ has
$|h(x_i)-f(x_i)|\leq\tau/2+a\leq5\tau/8$, so all training inputs
avoid $A_h$. Its test errors lie in $A_h$, since
$\tau-a\geq7\tau/8>3\tau/4$. A union bound excludes every surviving
center with $D(A_h)>\varepsilon/3$ at the displayed sample size.
The cover is common to all widths, so the assertion is simultaneous
without a union bound over widths. An exact teacher clone is feasible
with energy at most $E_*$, so every energy-minimizing exact interpolant
belongs to the covered sublevel. The risk bound persists under mixtures.
\end{proof}

\begin{corollary}[An exact Bayesian strong-prior limit]\label{cor:exact}
Use Gaussian coordinates in~\eqref{eq:strong-scales}. At every fixed
$N\geq C$, every iterated weak limit of the posterior obtained by first
sending $\sigma\downarrow0$ at fixed $\rho>0$, and then sending
$\rho\downarrow0$, is supported on minimum-energy exact interpolants.
Such subsequential limits exist. With the sample bound of
Proposition~\ref{prop:exact}, all of them obey the stated regression and
posterior-confidence guarantees.
\end{corollary}
\begin{proof}
Fix a clean sample and put $e_S(w)=\max_i|z_w(x_i)-f(x_i)|$,
$O_\rho=\{e_S<\rho\}$,
$b_\rho=\inf_{O_\rho}\kappa_N$, and
$b_0=\min_{\{e_S=0\}}\kappa_N$.
The exact feasible set is nonempty and closed. Since $\kappa_N$ is
continuous and coercive, its minimum exists, with $b_0\leq E_*$.

The Gaussian density is proportional to
$e^{-\kappa_N/(2\sigma^2)}$. For every fixed $\rho,\eta>0$, an
open ball inside $O_\rho$ has energy below $b_\rho+\eta/2$.
Comparing its Gaussian mass to the tail
$\{\kappa_N>b_\rho+\eta\}$ proves
\[
 \Pi_{\sigma,\rho}\{\kappa_N>b_\rho+\eta\}\longrightarrow0.
\]
This gives tightness, since energy sublevels are compact. Each posterior
has $e_S\leq\rho$ and $\kappa_N\geq b_\rho$ almost surely, so every
inner weak limit is supported on
$\{e_S\leq\rho,\ \kappa_N=b_\rho\}$.
This does not require identifying the closure of $O_\rho$ with the
whole non-strict feasible set.

As $\rho\downarrow0$, $b_\rho$ increases to $b_0$. Indeed choose
almost minimizers at tolerances tending to zero. They lie in a common
compact energy sublevel; a convergent subsequence has exact fit and
energy at most $\lim b_\rho$, giving the reverse inequality to
$b_\rho\leq b_0$.
The inner limiting measures are supported on the common compact set
$\{\kappa_N\leq b_0\}$, and every outer weak limit is supported on
$\{e_S=0,\kappa_N=b_0\}$. Proposition~\ref{prop:exact} then applies
directly to their support on its single high-probability sample event.
\end{proof}

The limit order in Corollary~\ref{cor:exact} is part of the statement.
It does not identify a fixed-scale regular conditional posterior given
exact outputs, reverse the two limits, or prove a result along arbitrary
joint paths $(\sigma,\rho)\to(0,0)$. No rank or regular-value assumption
on the sample-evaluation map was used. Under two-sided quadratic
tail envelopes with exponents $0<c_-\leq c_+$, the analogous safe energy
bound is $(c_+/c_-)E_*$, giving the same polynomial conclusion, without
identifying minimum-quadratic-energy interpolants.

\section{Binary classification and the role of a margin}\label{sec:classification}
The continuous-regression theorem does not assume that the teacher output
avoids zero. Classification requires a separate denominator argument.

\begin{corollary}[Classification on a support with a teacher margin]
\label{cor:classification}
Let $g:X\to\{-1,1\}$ and suppose a teacher of width $C$, depth $L$,
and weight bound $B$ satisfies $g(x)z_*(x)\geq r>0$ on $X$.
Condition the raw prior only on the training margins
$g(x_i)z_w(x_i)>1/2$. For both the strong-prior regime (at each fixed
$N\geq C$, for all $\sigma$ below a threshold that may depend on $N$, the
teacher, $r$, and $\varepsilon$, as in Theorem~\ref{thm:strong}) and every
fixed $\gamma\geq1$, there are fixed-depth sample bounds polynomial in
$d,C,B,1/r$ and the accuracy parameters, independent of learner width
for every $N\geq C$, such that with sample probability at least
$1-\delta$,
\[
 D\{x:\Pi(g(x)z_w(x)>0\mid S_n)\leq2/3\}\leq\varepsilon.
\]
The teacher margin may hold on any full-measure support; no assertion
about unobserved points outside that support is required.
\end{corollary}
\begin{proof}
Multiply the teacher readout by $3/(2r)$, obtaining signed output at
least $3/2$. Its width is still $C$, and its weight bound is at most
$\max\{B,3B/(2r)\}$. A uniform output tube of radius $1/2$ fits all
training margins. The preceding target-tube estimates therefore apply.
For the finite-cover step use radius $a<1/4$. If $h$ approximates a
fitted network, every training point satisfies $g(x_i)h(x_i)>1/2-a$.
Every sign error of that network lies in
$\{x:g(x)h(x)\leq1/2-a\}$, since $a<1/2-a$. A union bound over
these sets replaces Lemma~\ref{lem:posterior}; the small-width mass-only
argument is unchanged. The risk-to-confidence bound is again Markov.
\end{proof}

For $d\geq2$, the sphere $S^{d-1}$ is connected. A continuous network
with $|z_*|\geq r>0$ everywhere on it has a constant sign, by the
intermediate value theorem. Thus a nonconstant full-sphere binary target
does not satisfy the corollary's uniform-margin premise. The corollary is
useful, for example, for distributions whose support excludes a band around
a decision boundary. The disconnected sphere $S^0$ is an exception to
the connectedness observation.

For a nonconstant example on the entire connected sphere,
take the regression target $f(x)=\tanh(x_1)$, represented at depth two
and width one with weights bounded by one. Every $N\geq1$ is covered,
for every distribution on the sphere, and any fixed $0<\rho<\tau$.
This example needs no exclusion of an equatorial band. At larger fixed
depth, iterated compositions of tanh give corresponding width-one examples.

Real score labels also permit sign prediction without a uniform margin.
For any nonzero target value with $|f(x)|>\tau$, the event
$|z_w(x)-f(x)|\leq\tau$ implies the correct sign. Hence
\begin{equation}\label{eq:score-sign}
 D\{\Pi(\operatorname{sign}z_w=\operatorname{sign}f\mid S_n)\leq2/3\}
 \leq\varepsilon+D\{|f|\leq\tau\}
\end{equation}
on the regression theorem's success event. Any convention for the sign
of zero is harmless, since zeros lie in the second set. This is a
corollary for training on real-valued scores, not for training on binary
labels alone.

For $f(x)=\tanh(x_1)$ and uniform spherical inputs in dimension $d\geq2$,
\begin{equation}\label{eq:sphere}
 D\{|f|\leq\tau\}\leq3\sqrt d\,\tau\qquad(0<\tau\leq1/2).
\end{equation}
To check it, let $s=\operatorname{artanh}\tau$. Then
$s\leq4\tau/3$ and $s\leq\operatorname{artanh}(1/2)<0.55$.
For $d\geq3$, the first-coordinate density is
$c_d(1-t^2)^{(d-3)/2}$ with
$c_d=\Gamma(d/2)/(\sqrt\pi\Gamma((d-1)/2))\leq\sqrt d$;
the ratio bound follows by applying Jensen to the square root of a
Gamma variable. For $d=2$, the density
$1/(\pi\sqrt{1-t^2})$ is at most one on $[-s,s]$.
Integrating yields at most $2s\sqrt d\leq3\sqrt d\tau$.
Taking regression accuracy $\varepsilon/2$,
$\tau=\varepsilon/(6\sqrt d)$, and $\rho=\tau/2$ gives correct-sign
posterior confidence on all but an $\varepsilon$ fraction of the entire
sphere, with a fixed-depth sample bound polynomial in $d,1/\varepsilon$,
at every learner width $N\geq1$. The observed labels here remain the
real scores $\tanh(x_1)$.

\section{What this result does and does not establish}
The argument establishes a continuous-input version of width-independent
learning with a positive training tolerance. Its sample bound is polynomial
in teacher width and weight bound at fixed depth, for every learner width
above the teacher width. The distribution may have atoms, a density, or
lower-dimensional support; no covering-number assumption on the distribution
itself is made. The unit-ball assumption is used to bound the logarithm of
an input-net size and the dependence of outputs on parameters. A bounded
domain of a different radius can be rescaled, changing the teacher's
first-layer weight bound.

The positive tolerance in the finite-scale theorem is substantive. It does
not prove exact real-valued interpolation under a regular conditional
posterior or a zero-tolerance limit at fixed prior scale.
Section~\ref{sec:exact} supplies a distinct exact-interpolation result
with a specified order of limits. In the strong-prior tolerance theorem,
the sample polynomial can be independent of the absolute training tolerance,
but the scale threshold may deteriorate as that tolerance shrinks.
At finite scale the displayed bound tracks this cost through
$\log^+(1/\rho)$. The inverse gap $1/(\tau-\rho)$ also enters the
polynomial, so the theorem does not make a uniform assertion as the test
and training thresholds coincide.

The binary corollary assumes a uniform teacher margin on the data support.
Dropping that assumption would require a different lower bound on posterior
normalization, or quantitative control of how much input mass lies near the
teacher's decision boundary. The present proofs do not establish a
width-only, margin-free classification theorem. They also do not establish
the bounds for ReLU, noisy labels, $\gamma<1$, or efficient posterior
computation. No robustness restriction on the network prior is imposed,
and no polynomial dependence on depth is claimed.

\begingroup\small
\begin{thebibliography}{9}
\bibitem{nongaus} K. H\"anni, D. Vaintrob, and OpenAI Codex.
\emph{Bayesian neural-network learning with non-Gaussian weight priors}.
Research draft, 2026.
\url{https://github.com/kaarelh/AI-works/tree/main/neural-network-generalization/side-results/non-gaussian-priors}.
\bibitem{hoeffding} W. Hoeffding.
Probability inequalities for sums of bounded random variables.
\emph{Journal of the American Statistical Association} 58 (1963), 13--30.
\bibitem{brascamp-lieb} H. J. Brascamp and E. H. Lieb.
On extensions of the Brunn--Minkowski and Pr\'ekopa--Leindler theorems,
including inequalities for log concave functions, and with an application
to the diffusion equation.
\emph{Journal of Functional Analysis} 22 (1976), 366--389.
\bibitem{caffarelli} L. A. Caffarelli.
Monotonicity properties of optimal transportation and the FKG and related
inequalities. \emph{Communications in Mathematical Physics} 214 (2000), 547--563.
\bibitem{ledoux} M. Ledoux.
\emph{The Concentration of Measure Phenomenon}. AMS, 2001.
\end{thebibliography}\endgroup
\end{document}
